\section{Assumption 2}\label{s:ass2}

\lem{3}{There exist projects of objective worth: projects whose value does not depend on the attitudes of those engaged in them \citep{Metz2013}.}

\paragraph{The sentence as a whole.} Assumption~2 has two parts separated by a colon. The first part is an existential claim: there are projects of objective worth. The second part says what objective worth is: a project has it if its value does not depend on the attitudes of those engaged in it. The second part is grammatically an apposition, but it functions as a definition, and it is the only place in the text where \enquote{objective} is explained. Definition~1, Assumption~3 and the proof of Theorem~4 all use the phrase \enquote{of objective worth}, and all of them inherit its sense from the colon of Assumption~2.

\paragraph{The role of the assumption.} One would expect an assumption in an axiomatic text to be used in the proof of its theorem. Assumption~2 is not.

\begin{proposition}[Redundancy in the derivation]\label{p:redundant}
Theorem~4 follows from Definition~1, Assumption~3 and the step from ability to an available life (Proposition~\ref{p:bridge}), without Assumption~2. Moreover, Assumption~3 entails the existential part of Assumption~2 provided that there is at least one agent capable of intentional action.
\end{proposition}

\begin{proof}
The proof of Theorem~4 (\Lm{6}--\Lm{8}) cites Assumption~3 and Definition~1 and nothing else, and each of its steps is licensed by one of these or by the step from ability to an available life. For the second claim: if there is an agent capable of intentional action, Assumption~3 gives a project of objective worth in which she is able to engage, and so there is a project of objective worth.
\end{proof}

The Lean module confirms the first claim mechanically: Assumption~2 is not among the axioms on which \texttt{meaning\_available} depends (Section~\ref{s:lean}). The second claim shows that the existential part of Assumption~2 adds nothing to Assumption~3 in any world that contains an agent.

Why, then, is the assumption stated? Proposition~\ref{p:redundant} answers the question by elimination. The assumption does three things, none of them in the derivation.

First, it \emph{interprets}. Its second part fixes the sense of \enquote{objective worth}, and without it Definition~1 would be indeterminate between the hybrid view and the subjective view: \enquote{objective} might mean no more than \enquote{real} or \enquote{genuine}, and a subjectivist could accept every sentence of the text. The formal derivation goes through under any interpretation of the predicate (Proposition~\ref{p:neutralfml}); the philosophical content of Theorem~4 depends on its interpretation, and Assumption~2 supplies it.

Second, it \emph{separates}. The existential part isolates, as a claim on its own, what the hybrid view needs from the theory of value: that there is something with worth of the required kind. A reader who doubts this has a reason to doubt the text that is independent of any doubt about Assumption~3, and the text lets her locate her doubt. The separation also lets the reader see that Assumption~3 is stronger than Assumption~2, since it entails it given an agent and adds a claim about every agent's abilities.

Third, it \emph{answers}. Remark~5 says that Assumption~2 denies the premise of the sceptical challenge that, \emph{sub specie aeternitatis}, nothing matters (\Lm{11.4}). The assumption is stated so that the text can say exactly where it stands on that challenge: it does not refute it, but it is committed to denying one of its premises, and the commitment is made in a sentence of its own.

\lem{3.1}{There exist}

The existential quantifier ranges over projects. Two questions about its domain should be settled.

The first is whether projects are types or tokens. A token project is a particular person's undertaking: this mathematician's work on this problem, this nurse's care of these patients. A project type is something that many people can engage in: mathematical research, the care of the sick. Assumption~2 is best read as a claim about types, or about possible tokens, since it must support Assumption~3, which says that every agent has a project of objective worth in which she is \emph{able} to engage. A project in which no one has yet engaged can exist in the relevant sense. When a project type has objective worth, the token undertakings that instantiate it inherit its worth, though they may add or subtract worth according to how they are carried out.

The second question is whether the existence asserted is modest or robust. The modest reading says only that some projects meet the condition stated after the colon. The robust reading adds that objective worth is a property with a place in the ultimate inventory of the world, perhaps a non-natural one. The text is committed only to the modest reading. Whether worth is natural or non-natural, reducible or primitive, is a question for the theory of value; Assumption~2 says that, however that question is answered, some projects have value that does not depend on the attitudes of those engaged in them. ETH-01 and MTE-01 give a naturalist answer, which Section~\ref{s:defence} shows to be sufficient.

\lem{3.2}{projects of objective worth: projects whose value}

\paragraph{The colon.} The phrase repeats \enquote{projects}, which signals that what follows the colon specifies the class introduced before it. The definition is of the form: a project is of objective worth if and only if its value does not depend on the attitudes of those engaged in it. The definition presupposes that the project has a value, and makes objectivity a property of that value: objective worth is objective value.

\paragraph{Whose value?} The value in question is the value of the project, but a project's value usually derives from the value of its object. The value of caring for a sick person derives from the value of her well-being; the value of a mathematical inquiry derives from the value of understanding. EXI-01-c makes the same point in the proof of its Lemma~2.3: the value of caring for a person or of coming to understand something derives from the value of the person or of the understanding, not from the agent's attitude. This matters for how Assumption~2 is defended. It need not be shown that projects, as such, have intrinsic value. It need only be shown that their objects have value, or that the projects serve value, independently of the attitudes of those engaged. The value of a project of objective worth may be wholly derivative from the value of what it serves, and its worth is no less objective for that.

\lem{3.3}{does not depend on}

\enquote{Depend} admits two readings, and the choice between them determines what Assumption~2 claims.

On the \emph{counterfactual} reading, the value of a project does not depend on certain attitudes if it would have the same value whatever those attitudes were. On the \emph{grounding} reading, the value of a project does not depend on certain attitudes if those attitudes are not among the facts in virtue of which the project has its value.

The counterfactual reading is too strong. Consider a friendship. Its value depends, counterfactually, on the attitudes of the friends: if they did not care about one another, there would be no friendship and so no value of friendship. On the counterfactual reading, friendship would therefore lack objective worth, and Definition~1 would imply that friendship confers no meaning. This is absurd, and it cannot be what the text means, since the family counts relationships of care among the paradigms of meaningful engagement (EXI-01-c, its Assumption~2.2; EXI-17-a, its Definition~2.2).

The grounding reading avoids this result if \enquote{attitudes} is read with care. The friends' care for one another is a constituent of the friendship, not an evaluation of it. What Assumption~2 excludes is that the value of the friendship should be grounded in the friends' \emph{evaluative} attitudes toward it: in their regarding it as valuable, or their being attached to it, or their caring about it as a project. A friendship is not valuable because the friends value it; they value it, if they see correctly, because it is valuable, and it is valuable in part because it consists in mutual care.

\begin{proposition}[Restricted objectivity]\label{p:restricted}
On the reading required by the text, a project is of objective worth if and only if it has value, and that value is not grounded in the evaluative attitudes, toward the project or its object, of those engaged in it.
\end{proposition}

\begin{proof}[Argument]
The counterfactual reading of \enquote{depend} implies that relationships of mutual care have no objective worth, which contradicts the family's treatment of such relationships as paradigms of meaningful engagement; so the text requires the grounding reading. On the grounding reading, if \enquote{attitudes} included every attitude, the same contradiction would follow, since mutual care consists of attitudes; so the attitudes excluded are evaluative attitudes, those by which the engaged agents regard, prize or care about the project or its object as worth their engagement. The resulting condition is what Definition~1 needs: it excludes the case in which engagement confers on its object the worth that makes the engagement meaningful.
\end{proof}

Proposition~\ref{p:restricted} identifies the target of Assumption~2 precisely. The target is the view that engagement \emph{bootstraps}: that by caring about something, one makes it worth caring about, so that every caring engagement is engagement with something of worth. That view is Taylor's, and in a qualified form Frankfurt's. Assumption~2 denies that bootstrapping is the only source of worth: it says that there are projects whose worth is not bootstrapped.

\lem{3.4}{the attitudes of those engaged in them}

\paragraph{The restriction.} Assumption~2 does not say that the value of a project of objective worth is independent of all attitudes. It says that it is independent of the attitudes of \emph{those engaged in it}. This is a restriction, and it is easy to miss.

The restricted notion is compatible with three kinds of attitude-dependence that a stronger notion would exclude. It is compatible, first, with the dependence of a project's value on the attitudes of those it serves: the value of giving a gift depends in part on the pleasure of the recipient, and the value of teaching depends in part on the pupils' coming to understand, which is an attitude. It is compatible, second, with the dependence of value on the attitudes of a community, as the value of preserving a language may depend on the attachment of its speakers to it. It is compatible, third, with theories on which value in general is constituted by idealized attitudes: theories on which something is good if and only if idealized agents would desire or endorse it. On such theories, the value of a project depends on attitudes, but not on the actual attitudes of the people engaged in it, which may differ from the idealized ones.

\paragraph{Rival readings in the collection.} The manuscripts that restate Assumption~2 do not all preserve the restriction. EXI-17-a defines a worthwhile object as one that has or serves value that does not depend on \enquote{the agent's attitudes toward it} (its Definition~2.2), which is the restricted notion, made more precise by the phrase \enquote{toward it}, which agrees with Proposition~\ref{p:restricted}. EXI-01-c adopts the same definition (its Definition~2.1). EXI-06-b, by contrast, describes meaning on the account of EXI-01 as engagement with what has worth \enquote{independently of anyone's attitudes} (see the testimonia, Appendix~\ref{app:test}). That is the unrestricted notion. EXI-06-b's argument does not need it: the argument is that the universe's lack of attitudes toward us takes nothing from the worth of what we engage with, and on the restricted notion the universe is not among those engaged in our projects, so its attitudes are irrelevant in any case. EXI-06-b's gloss is therefore a strengthening of the text that its argument does not require, and the Commentary reads Assumption~2 as restricted.

\begin{proposition}[Sufficiency of the restricted notion]\label{p:suff}
Every use that the collection makes of Assumption~2 is supported by the restricted notion of objective worth.
\end{proposition}

\begin{proof}[Argument]
The uses are four. (i) Definition~1 needs objective worth to exclude engagement with worthless objects such as Taylor's stones and Rawls's grass; the restricted notion does this, since the only source of worth for those objects would be the agent's own attitudes. (ii) EXI-17-a's argument that meaning is not hostage to luck in outcomes (its Corollary~3.6) needs only that the worth of the object belong to the object rather than to the agent's success, which the restricted notion secures. (iii) EXI-06-b needs only that worth not depend on the universe's attitudes, and the universe is not engaged in our projects. (iv) EXI-04-a needs, in its reply to its Objection~2, only that worth not be conferred by the seriousness with which we pursue our ends, and seriousness is an attitude of the engaged agent. The theistic case is the one that might seem to need more: EXI-11-a holds that worth is not conferred by assignment (its Principle~P1), so that a creator's purposes for us do not make our activities worthwhile. But Principle~P1 of EXI-11-a is a separate principle, argued for separately, and EXI-01-a does not need it to state its account.
\end{proof}

\paragraph{\enquote{those engaged}.} The plural noun phrase covers projects in which many people are engaged together. A scientific research program, an orchestra, a political movement are projects with many participants, and Assumption~2 requires that their value not be grounded in the evaluative attitudes of any of them. This excludes a form of collective bootstrapping, on which a shared enthusiasm makes its object worthwhile for those who share it. A community of enthusiasts for a worthless pursuit is not given meaning by its enthusiasm, however widely shared, any more than a single enthusiast is.

\lem{3.5}{\citep{Metz2013}}

The citation of Metz rather than Wolf at this point is significant. Wolf supplies the account of meaning in Definition~1, but she declines to give a theory of objective worth. Metz gives the most sustained recent argument that meaning depends on conditions that are not constituted by the subject's attitudes. His survey of theories argues, against subjectivism, that it cannot account for the judgment that some engaged lives lack meaning, and, against supernaturalism, that meaning does not require a relation to a spiritual being or realm; his own fundamentality theory holds that a life is more meaningful the more its subject employs her reason in positive ways toward conditions of human existence that are responsible for much else about it \citep{Metz2013}.

What EXI-01-a takes from Metz is the existential claim and the argument against subjectivism, not the fundamentality theory. The fundamentality theory is a rival to Definition~1 as an account of what meaning consists in: it makes the degree of meaning depend on the fundamentality of the conditions to which one's reason is directed, a notion that Definition~1 does not use. The citation is therefore a citation of an ally on one point, the existence of attitude-independent worth, and not of a source of the definition. A reader who checks the citation will find in Metz a theory that disagrees with EXI-01-a about the measure of meaning while agreeing with it about the existence of what meaning is a measure of.

\subsection{A defence of Assumption 2}\label{s:defence}

Remark~5 says that a defence of Assumption~2 is given in the Commentary (\Lm{12}). We give it here. It establishes the restricted claim of Proposition~\ref{p:restricted}: that there are projects whose value is not grounded in the evaluative attitudes of those engaged in them. It does not establish any stronger claim about the metaphysics of value, and it does not answer the demand, considered in Section~\ref{s:absurd}, that worth be vindicated from a standpoint that presupposes nothing. The defence has four parts.

\paragraph{(D1) The argument from cases.} Taylor's Sisyphus, given the desire to roll stones, is actively engaged in rolling them. Rawls's grass counter is actively engaged in counting. If the value of a project were grounded in the evaluative attitudes of those engaged in it, these projects would have as much value as any, since their subjects prize them as much as anyone prizes anything, and engagement with them would confer as much meaning as engagement with anything. Reflection judges that it does not: the lives of the stone-roller and the grass counter are not meaningful in respect of their rolling and counting, however engaged they are, and they are less meaningful than the lives of people equally engaged in the care of others or in the pursuit of understanding. The best explanation of this judgment is that the value of the care of others and of understanding is not grounded in the attitudes of those engaged in them, and the value of rolling stones and counting grass, if any, is. So there are projects of objective worth in the restricted sense.

The subjectivist can reject the judgment. Taylor did: he held that his Sisyphus's life was meaningful. But the cost of rejecting it is high, since it requires holding that there is no difference in meaning between the life of the grass counter and the life of an equally engaged physician, and few who consider the case are prepared to say this. Judgments about cases are evidence in this domain, as EXI-17-a assumes (its Assumption~2.4) and as ETH-01 holds for judgments about the good. The subjectivist's position is consistent, but it is not supported by the evidence that bears on the question.

\paragraph{(D2) The argument from the content of engagement.} \citet{Wiggins1976} observed that a purely subjective account of meaning cannot be held by someone who finds her life meaningful, because engagement presents its object as worth engaging with in a way that does not depend on the engagement. The mathematician does not care about her problem and conclude that, since she cares about it, it is important. She cares about it because she takes it to be important, and if she came to believe that its importance was nothing but a projection of her caring, her caring would be undermined. The same is true of the nurse's care for her patients and of the parent's love for a child: in each case the engaged agent takes the worth of the object to be the reason for her engagement and not its product.

If subjectivism were true, engagement would therefore be systematically mistaken about its own content: every engaged agent would take her object to have a worth it has only in virtue of her taking it to have it. This would make every meaningful life rest on an illusion. EXI-18 holds that some illusions are psychologically necessary, but that their necessity does not make them true, and EXI-18-a's reply to the objection that meaning needs illusion is that the serious agent need not believe that her ends are certified from the standpoint of the universe, only that they are worthwhile, which may be true (its reply to Objection~3). Subjectivism denies that it may be true in the sense in which the agent believes it. The hybrid view, and Assumption~2 with it, allows the engaged agent's view of her object to be correct.

\paragraph{(D3) The argument from error and discovery.} People discover that their lives have been wasted. Tolstoy's Ivan Ilyich, in \emph{The Death of Ivan Ilyich}, comes to see on his deathbed that the life he had lived, correct, successful and absorbed in its own advancement, had not been what it should have been. Mill records in his \emph{Autobiography} (ch.~V) how, at twenty, he asked himself whether the realization of all the reforms to which he had devoted himself would make him happy, answered that it would not, and found that the ends that had given his life its direction had lost their hold on him. Tolstoy's \emph{Confession} describes a crisis of the same structure \citep{Tolstoy1987}. In each case the subject had been engaged; in Ivan Ilyich's case the discovery is that the objects of his engagement had not deserved it.

On a subjectivist account, such discoveries are unintelligible as discoveries. If the worth of Ivan Ilyich's pursuits was grounded in his attitudes toward them, then while he held those attitudes his pursuits were worthwhile, and his deathbed judgment can only record a change of attitude, not a recognition that he had been wrong. But the judgment presents itself as a recognition, and it is so understood by those who read the story. The possibility of being wrong about whether one's life is meaningful, and of discovering it, requires that meaning depend on something that one's attitudes do not settle. Assumption~2 says what that something is.

Mill's case shows a different structure, and it is important not to assimilate it. Mill did not discover that the reforms to which he had devoted himself were worthless; he never doubted their worth. He discovered that his engagement had failed, that the ends no longer moved him. On Definition~1, his life during the crisis lacked meaning not because its objects lacked worth but because it lacked engagement, and his recovery, which he attributed to the cultivation of feeling, restored the subjective component. The pair of cases illustrates the two components of Definition~1: Ivan Ilyich had engagement without worth, Mill had worth without engagement.

\paragraph{(D4) The argument from the theory of value.} The collection contains a theory of value that entails Assumption~2. ETH-01 holds that the good is what idealized reflective endorsement converges on: one outcome is better than another if and only if every agent's idealized reflective endorsement would prefer the first (Definition~3.2 of ETH-01-a), where idealized reflective endorsement is the attitude an agent would take if she were fully informed, vividly imagined the alternatives, reasoned without error and gave equal weight to everyone affected (Definition~3.1 of ETH-01-a). MTE-01 holds that moral facts are natural facts about idealized reactions.

\begin{proposition}[ETH-01 entails Assumption~2]\label{p:eth01}
If the good is what idealized reflective endorsement converges on, and the good has the components listed in ETH-01, then there are projects of objective worth in the restricted sense.
\end{proposition}

\begin{proof}
Take a project that serves one of the components, for instance the care of the sick, which serves their well-being. On ETH-01, the outcomes this project brings about are better than their alternatives because the idealized reflective endorsement of every agent would prefer them. That fact does not depend on the actual evaluative attitudes of those engaged in the project: had the nurses who care for the sick regarded the care as worthless, or been indifferent to it, the idealized endorsement of every agent, theirs included, would still prefer the outcomes in which the sick are cared for, since idealization corrects for ignorance, failures of imagination and partiality, and nothing in the nurses' actual attitudes enters the definition except as material for correction. So the value of the project is not grounded in the evaluative attitudes of those engaged in it, and it is of objective worth in the sense of Proposition~\ref{p:restricted}.
\end{proof}

Proposition~\ref{p:eth01} shows that Assumption~2 is no more doubtful than the theory of the good of ETH-01, and it shows why the restriction in \Lm{3.4} matters. On ETH-01, value is attitude-dependent in a sense: it is constituted by idealized attitudes. An unrestricted notion of objective worth, on which worth is independent of all attitudes, would be inconsistent with ETH-01. The restricted notion is not. The text's restriction is therefore what allows EXI-01-a to rest on the theory of value of the collection rather than on a more contested non-naturalism.

The standard objections to that theory have answers elsewhere in the collection, which we record by their conclusions. The argument from queerness, that objective values would be entities of a strange sort known by a strange faculty \citep{Mackie1977}, proves too much: it would make mathematics false as well (MTE-05). The argument from evolutionary debunking, that our evaluative dispositions were selected for reasons unrelated to their truth \citep{Street2006}, fails because evolution selected for tracking reasons, and those reasons include evaluative ones (MTE-09, stated for moral reasons). The argument from disagreement fails because disagreement about value is no deeper than disagreement in other theoretical domains that are not thought to lack facts (MTE-12). Each of these is stated in the collection for moral facts, and each applies to the restricted worth of Assumption~2 at least as well, since the restricted claim is weaker than moral realism.

\paragraph{Objections to the defence.}

\begin{objection}[Frankfurt]\label{o:frankfurt}
\citet{Frankfurt1982,Frankfurt2004} argues that what is important to a person is important to her largely because she cares about it, and that love, in particular, does not respond to the value of its object but confers value on it for the lover. The parent does not love her child because the child is valuable; the child is valuable to her because she loves it. So the paradigm of meaningful engagement is a case in which worth is conferred by the attitude of the one engaged.
\end{objection}

\begin{reply}
Frankfurt's claim concerns importance \emph{to} the lover, and on that it may be right: the child matters to the parent in a way it matters to no one else, and the difference is made by her love. Assumption~2 concerns the worth of the project, here the care of the child, and that worth is not conferred by the parent's love. The child's well-being has value whether or not anyone loves her, as the case of an unloved child shows: it is not a matter of indifference whether an unloved child flourishes. A parent who cares for her child is engaged with something whose worth she did not create, and her love is fitting because of that worth. What her love adds is a particular relation to that worth, the relation that EXI-16 calls a reason of love not transferable to qualitatively similar persons; it does not add the worth. Frankfurt's view of importance can be accepted as an account of the subjective component of Definition~1, which is what the comment on \Lm{2.5} took from him, and rejected as an account of the objective component.
\end{reply}

\begin{objection}[Elitism]\label{o:elitism}
Every attempt to say which projects have objective worth privileges the projects of the educated and the privileged: art, science, causes. Assumption~2 therefore smuggles a class judgment into the theory of meaning.
\end{objection}

\begin{reply}
Assumption~2 is existential. It says that there are projects of objective worth and what their worth consists in; it does not say which projects have it. The list of goods that the collection uses, from ETH-01, begins with well-being and relationships, which are within the reach of every life, and EXI-01-d answers the objection directly: caring for others, honest work that serves others' needs and the pursuit of understanding at any level meet the condition of Assumption~2, and are open to everyone. The objection would hold against a theory that measured worth by prestige. It does not hold against a theory that measures it by idealized endorsement, which gives equal weight to everyone affected.
\end{reply}

\begin{objection}[Alienation]\label{o:alien}
If the worth of my projects is independent of my attitudes, it is not clear why engaging with them makes \emph{my life} meaningful rather than merely making the world better. Objective worth, so understood, is alien to me.
\end{objection}

\begin{reply}
The objection would hold against a purely objective theory, on which meaning consists in serving objective value whether or not one cares about it. On such a theory, the meaning of my life would be a fact about the world's ledger. Definition~1 is not such a theory. It makes meaning consist in a relation between the subject and objective worth, the relation of active engagement, and the relation is as much the subject's as anything in her life. Objective worth is what makes the relation fitting; engagement is what makes it hers. This is the reason Wolf gave for the hybrid view, and it is the reason Definition~1 has two components.
\end{reply}

\paragraph{What the defence establishes.} The defence establishes that Assumption~2 is supported by judgments about cases (D1), by the content of engagement itself (D2), by the intelligibility of discovering that one's life has been wasted (D3), and by the theory of value of the collection (D4), and that the objections from Frankfurt, from elitism and from alienation do not undermine it. It does not establish that the theory of value is true; it shows that Assumption~2 is supported to the extent that that theory is, and that the restricted form of the assumption is weaker than the theory. Nor does it answer the sceptic who grants all this and asks what vindicates the judgments about cases, the content of engagement and the theory of value from a standpoint that does not presuppose them. That sceptic is Nagel's, and Remark~5 says that Assumption~2 denies the premise of his challenge rather than refuting it. Section~\ref{s:absurd} examines what this means. Volume~II, Part~A, develops (D4) in detail, and Part~B examines the subjectivist and hybrid theories at length.

\subsection{Assumption 2 and the metaethical options}\label{s:metaeth}

The defence has shown that Assumption~2 follows from the theory of value of ETH-01. A reader who does not accept that theory will want to know whether Assumption~2 survives under the others. Because the assumption is restricted (Proposition~\ref{p:restricted}), it is compatible with more metaethical views than an unrestricted assumption of objective worth would be. We take the main options in turn, describing each by its central thesis and, where the collection has a family on it, by that family's conclusion.

\paragraph{Non-naturalist realism.} On the view that evaluative properties are sui generis and not natural, worth is independent of every attitude, and Assumption~2 holds in its unrestricted form and so in its restricted form. The restriction costs the non-naturalist nothing.

\paragraph{Naturalist realism.} On the view of MTE-01, moral facts are natural facts about idealized reactions; on the view of ETH-01, the good is what idealized reflective endorsement converges on. Proposition~\ref{p:eth01} has shown that Assumption~2 holds on these views because the idealized attitudes that constitute value are not the actual attitudes of those engaged. The restriction is what makes Assumption~2 consistent with them: an unrestricted assumption would deny that value is constituted by attitudes of any kind.

\paragraph{Subjective response-dependence.} On the view that something is valuable for a person if and only if she would desire it under some modest idealization of her own attitudes, value is relative to the person. Assumption~2 is false on this view if the idealization is so modest that the grass counter, idealized, would still desire to count grass; it is true if the idealization corrects for the narrowness that makes the grass counter's desire compulsive. The defence of (D1) argues that the case requires the second. A response-dependent view that cannot distinguish the grass counter from the physician is not a rival to Assumption~2 but a version of subjectivism, and is answered by (D1) to (D3).

\paragraph{Expressivism and quasi-realism.} On expressivist views, evaluative judgments express non-cognitive states; MTE-03 holds that moral sentences express states of being for, which compose according to the logic of plans. Such views have a characteristic way with claims of attitude-independence: the claim that the value of caring for the sick does not depend on the attitudes of the carers is itself an evaluative claim, expressing a plan to value the care whatever the carers' attitudes, and the expressivist can make it. Assumption~2 can therefore be asserted by an expressivist, as a first-order claim. MTE-17 adds that quasi-realism, once its explanatory commitments are made explicit, is indistinguishable from realism, and on that conclusion the expressivist who asserts Assumption~2 asserts what the realist asserts. Either way, the assumption is not hostage to the outcome of the dispute between realism and expressivism.

\paragraph{Constructivism.} On constructivist views, evaluative truths are the outcomes of an idealized procedure of practical reasoning; MTE-18 holds that constructivism is realist about procedures and not about their outcomes, and that the procedure itself is mandatory. On such views the worth of a project is the verdict the procedure delivers on it. The procedure is not the actual attitudes of those engaged in the project, and Assumption~2 in its restricted form holds whenever the procedure delivers a verdict of worth on some project independently of what its participants happen to think, which any procedure worth the name does.

\paragraph{Error theory.} On the view that all evaluative judgments are false, because they presuppose objective values that do not exist \citep{Mackie1977}, Assumption~2 is false. MTE-05 answers this view: the argument from queerness proves too much, since it would make mathematics false as well. The error theorist is the one metaethical opponent of Assumption~2, and the dispute with her is not special to the theory of meaning.

\begin{proposition}[Metaethical robustness]\label{p:robust}
Assumption~2 in its restricted form holds on every metaethical view considered above except error theory and those forms of subjective response-dependence that cannot distinguish idealized from actual desire.
\end{proposition}

\begin{proof}[Argument]
By the discussion of each view: non-naturalist realism entails the unrestricted form; naturalist realism of the kind of ETH-01 and MTE-01 entails the restricted form (Proposition~\ref{p:eth01}); expressivism permits it as a first-order claim and, on MTE-17, does not differ from realism in asserting it; constructivism entails it whenever its procedure is not identical with the actual attitudes of participants; error theory denies it; and response-dependence entails it just when its idealization corrects for actual desire.
\end{proof}

Proposition~\ref{p:robust} explains why EXI-01-a can state Assumption~2 without defending a metaethics. The assumption is a point on which almost every metaethical view that allows evaluative truth agrees, once the restriction of \Lm{3.4} is observed. Its opponents are the error theorist, who denies that there is evaluative truth, and the subjectivist about meaning, who holds that the only evaluative truth relevant to meaning is truth about what the engaged agent herself wants. The first is answered by MTE-05 and the second by the defence.
